% Figure: organic Rankine cycle on the T-s plane, drawn schematically against
% the saturation dome of a dry working fluid (a fluid whose saturated-vapour
% line leans to the right, so expansion from saturated vapour ends dry). The
% cycle sits at low temperature, recovering heat from an industrial exhaust
% stream that the dashed line marks. The contrast with the water Rankine dome
% is the rightward lean of the vapour line, which is why the turbine exit needs
% no superheat to stay out of the wet region.
\begin{figure}[htbp]
\centering
\begin{tikzpicture}
\begin{axis}[
  width=12cm, height=8.5cm,
  xlabel={specific entropy $s$ (kJ/kg$\cdot$K)},
  ylabel={temperature $T$ ($^{\circ}$C)},
  xmin=0.8, xmax=2.4, ymin=0, ymax=200,
  axis lines=left,
  tick label style={font=\scriptsize}, label style={font=\small},
  clip=false,
  ytick={0,50,100,150,200},
]
% === saturation dome for a dry fluid: the vapour line leans right ===
\addplot[enggray, thick, smooth] coordinates {
  (1.05,5) (1.10,40) (1.16,75) (1.22,105) (1.28,125) (1.34,135)
};
\addplot[enggray, thick, smooth] coordinates {
  (1.34,135) (1.55,128) (1.78,112) (1.98,92) (2.12,65) (2.20,30)
};
\node[font=\scriptsize, enggray, anchor=south] at (axis cs:1.34,135)
  {critical point};
% === the industrial waste-heat source temperature (dashed) ===
\addplot[engamber, thick, dashed, samples=2, domain=0.8:2.4] {150};
\node[font=\scriptsize, engamber, anchor=south west] at (axis cs:0.82,150)
  {exhaust source $\approx 150\,^{\circ}$C};
% === the cycle states ===
% 1 pump inlet (sat liquid at condenser, ~35 C), 2 pump exit (~36 C),
% 3 evaporator exit / turbine inlet (sat or slightly superheated vapour ~120 C),
% 4 turbine exit (condenser pressure, dry vapour ~40 C).
\addplot[engblue, very thick] coordinates {
  (1.10,35) (1.12,37)
};
% 2 to 3 along the high-pressure isobar: liquid heat, then through the dome
\addplot[engred, very thick] coordinates {
  (1.12,37) (1.20,90) (1.27,120) (1.45,122)
};
% 3 to 4 turbine expansion (nearly isentropic, ends dry to the right of dome)
\addplot[engblue, very thick] coordinates {
  (1.45,122) (1.50,40)
};
% 4 to 1 condenser (constant temperature, leftward)
\addplot[engblue, very thick] coordinates {
  (1.50,40) (1.10,35)
};
% labels
\node[font=\scriptsize, anchor=east] at (axis cs:1.10,35) {1,2};
\node[font=\scriptsize, anchor=south] at (axis cs:1.45,122) {3};
\node[font=\scriptsize, anchor=west] at (axis cs:1.50,40) {4};
\node[font=\small, engred, anchor=west] at (axis cs:1.25,100)
  {evaporator};
\node[font=\small, engblue, anchor=west] at (axis cs:1.52,80)
  {turbine};
\end{axis}
\end{tikzpicture}
\caption[Organic Rankine cycle on the temperature-entropy plane]{An organic
Rankine cycle on the $T$-$s$ plane, drawn against the saturation dome of a dry
working fluid (grey). The fluid is pumped (1 to 2), heated and evaporated by
the industrial exhaust stream along the high-pressure isobar (2 to 3), expanded
through the turbine (3 to 4), and condensed back to saturated liquid (4 to 1).
The vapour line leans to the right, the mark of a dry fluid, so the expansion
from state 3 ends to the right of the dome at state 4 and the turbine exit stays
dry without any superheat, which protects the blades and removes the need for the
high superheat a water cycle requires. The whole cycle sits below the exhaust
source temperature (amber, dashed) and well below the temperatures of a
steam Rankine plant, which is what lets it harvest heat that a steam cycle
could not economically raise into work.}
\label{fig:vol04ch03:orc-ts}
\end{figure}
